%% =====================================================================
%%  M.S. THESIS  --  Achuthan Rathinam  --  Boston University, CDS
%%
%%  What Does a Physics Network Actually Learn?
%%  Mechanistic Interpretability Across a Spectrum of
%%  Physics-Informed Neural Architectures
%%
%%  Built on the BU CDS Masters Thesis Template (buthesis.cls):
%%  https://github.com/BU-CDS/cds-masters-thesis-template
%%
%%  Build:   latexmk -pdf thesis.tex          (from docs/thesis/)
%%  Clean:   latexmk -c
%%
%%  Class options:
%%     notes    (default) show the \note{...} / \todo{...} draft notes
%%     nonotes  hide them all -- use for the final PDF
%%
%%  This file is currently an OUTLINE: each chapter has a one- or
%%  two-line note on what it will cover plus its section headings.
%%
%% ---------------------------------------------------------------------
%%  DECISIONS TO LOCK BEFORE WRITING
%%  * RQ5 (grokking) is exploratory; Ch. 8 is sized to shrink if thin.
%%  * Degree name: "Master of Science" here; the template's default is
%%    "Masters of Data Science".  Confirm with Dan Yang (CDS Graduate
%%    Affairs, danybu@bu.edu).
%%  * Readers' official titles.  No "Dr." before names.
%%  * Bibliography style: numbered by order of appearance (unsrtnat, as
%%    in the proposal) or alphabetical (plainnat).  Never mix.
%%  * Graduation year: set to 2026.  Must be the year of official
%%    conferral, not submission -- change to 2027 if conferral is in
%%    January.
%% =====================================================================
\PassOptionsToPackage{numbers,sort&compress}{natbib}   % numbered citations
\documentclass[notes]{buthesis}

\usepackage{amsmath}
\usepackage{booktabs}
\usepackage{xcolor}
\setcitestyle{square}
\emergencystretch=1in   % let long contents entries wrap instead of running into the right margin

% ---------- draft-note macros (hidden by the nonotes class option) ----
\newcommand{\note}[1]{\templatenote{{\color{blue!55!black}\sffamily [#1]}}}
\newcommand{\todo}[1]{\templatenote{{\color{red!70!black}\sffamily\bfseries [TODO: #1]}}}

% ---------------------------------------------------------------------
%  Your information
% ---------------------------------------------------------------------
\school{FACULTY OF COMPUTING AND DATA SCIENCE}
\doctype{Thesis}
% Title page: all capitals, bold, centered. Break lines with \\
\title{WHAT DOES A PHYSICS NETWORK ACTUALLY LEARN?\\
       MECHANISTIC INTERPRETABILITY ACROSS A\\
       SPECTRUM OF PHYSICS-INFORMED NEURAL\\
       ARCHITECTURES}
% The class spaces the title page for a two-line title; this one needs
% four, so take the extra lines out of the gap above "Submitted in...".
\makeatletter
\patchcmd{\maketitlepage}{\vspace{8\wordline}}{\vspace{4\wordline}}{}{\PatchFailed}
\makeatother
% Abstract page title (may be broken differently)
\abstracttitle{WHAT DOES A PHYSICS NETWORK ACTUALLY LEARN?\\
               MECHANISTIC INTERPRETABILITY ACROSS A\\
               SPECTRUM OF PHYSICS-INFORMED NEURAL\\
               ARCHITECTURES}
\author{ACHUTHAN RATHINAM}                            % title page
\authorfull{ACHUTHAN RATHINAM}                        % copyright page
\priordegrees{B.S., \todo{Institution}, \todo{Year}}  % use \\ for several
\degree{Master of Science}
\degreeyear{2026}

\begin{document}

% ------------------------- Preliminary pages -------------------------
\prelims                                  % roman numerals, footer
\maketitlepage                            % i   (no number shown)
\makecopyrightpage                        % ii  (no number shown)
% ---------------------------------------------------------------------
%  Approval page. Add or remove \reader lines as needed.
%  Third argument: title line(s); separate multiple lines with \\
%  The unsigned copy stays in the PDF; official titles; no "Dr.".
% ---------------------------------------------------------------------
\begin{approvalpage}
\reader{First Reader}{Siddharth Mishra-Sharma, Ph.D.}{Assistant Professor of Computing \& Data Sciences}   % TODO confirm official title
\reader{Second Reader}{Naomi Saphra, Ph.D.}{Assistant Professor of Computing \& Data Sciences}
\end{approvalpage}
              % iii (no number shown)
% ---------------------------------------------------------------------
%  Acknowledgments (optional)
% ---------------------------------------------------------------------
\buheading{ACKNOWLEDGMENTS}

\note{Advisor, readers, collaborators, compute resources, funding.}
       % optional
% ---------------------------------------------------------------------
%  Abstract (required). Title and name are taken from thesis.tex
%  (\abstracttitle and \authorfull). No figures or tables.
% ---------------------------------------------------------------------
\begin{buabstract}
\note{Write last. One paragraph each on the problem, the method, the
results, and the conclusions.}
\end{buabstract}
              % required
\tableofcontents                          % required
% ---------------------------------------------------------------------
%  List of Tables / Figures
%  Each list is REQUIRED if the thesis contains that kind of item
%  (main text or appendices). Delete a list if you have none.
% ---------------------------------------------------------------------
\listoftables
\note{Fills in automatically once tables exist.}

\listoffigures
\note{Fills in automatically once figures exist.}
                 % tables / figures
% ---------------------------------------------------------------------
%  List of Abbreviations (optional)
%  Entries MUST be in alphabetical order.
% ---------------------------------------------------------------------
\buheading{LIST OF ABBREVIATIONS}

\abbrev{CKA}{Centered Kernel Alignment}
\abbrev{EGNN}{E($n$) Equivariant Graph Neural Network}
\abbrev{GNN}{Graph Neural Network}
\abbrev{HNN}{Hamiltonian Neural Network}
\abbrev{MLP}{Multilayer Perceptron}
\abbrev{MSE}{Mean Squared Error}
\abbrev{RQ}{Research Question}
\abbrev{SO(3)}{Special Orthogonal Group in Three Dimensions (rotations)}
         % optional

% ---------------------------- Main text ------------------------------
\mainbody                                 % arabic numerals, header
\chapter{INTRODUCTION}
\label{ch:introduction}

\note{Motivates the problem, poses the central question, states the
hypothesis and the five research questions, and lists the intended
contributions. Adapted from the proposal's Problem Statement and Proposed
Research; Contributions and Organization are written last.}

\section{Motivation}
\section{The Central Question}
\section{Why Coulomb Dynamics}
\section{Thesis Statement}
\section{Research Questions}
\section{Contributions}
\section{Thesis Organization}

\chapter{BACKGROUND AND RELATED WORK}
\label{ch:background}

\note{Everything the reader needs to follow Chapters 3--8: the spectrum of
physics priors, each architecture family, the interpretability tools, and
the closest prior work (Cranmer et al.\ 2020 in particular). Related-work
survey still in progress.}

\section{Learning Physical Dynamics with Neural Networks}
\section{A Spectrum of Physics Priors}
\section{Symmetry-Constrained Architectures}
\subsection{Equivariance}
\subsection{Equivariant Graph Neural Networks}
\subsection{Attention-Based Equivariant Message Passing}
\subsection{Equivariant Transformers}
\section{Physics-Explicit Architectures}
\subsection{Hamiltonian Neural Networks}
\subsection{Hamiltonian Networks at Multi-Particle Scale}
\section{Unconstrained Baselines}
\section{Mechanistic Interpretability}
\subsection{Linear Probing and Its Pitfalls}
\subsection{Causal Intervention}
\subsection{Representational Similarity}
\subsection{Grokking and Delayed Generalization}
\section{Interpretability of Physics and Scientific ML Models}
\section{Summary and Gap}

\chapter{The Coulomb Dynamics Testbed}
\label{ch:testbed}

\note{Fully specifies the physical system, simulator, verification checks,
prediction task, dataset, and the analytic ground-truth quantities that
serve as probe targets (the ``answer key''). Decision pending: single-step
acceleration vs.\ multi-step position prediction.}

\section{Physical System}
\section{Simulator}
\section{Physics Verification}
\section{Prediction Task}
\section{Dataset}
\section{Ground-Truth Physical Quantities}

\chapter{Architectures and Training}
\label{ch:architectures}

\note{Defines the architecture suite, per-model implementation details
(including where ``layer $k$ activations'' are read from), the matched
training protocol, and shows comparable output accuracy across models so
that later representational differences are not explained by one model
being worse. Decision pending: five vs.\ six architectures.}

\section{Architecture Suite}
\section{Implementation Details}
\subsection{Multilayer Perceptron}
\subsection{Non-Equivariant Graph Neural Network}
\subsection{EGNN}
\subsection{Attention EGNN}
\subsection{Equivariant Transformer}
\subsection{Hamiltonian Neural Network}
\section{Matched Training Protocol}
\section{Output Accuracy}
\section{Equivariance Checks}

\chapter{INTERPRETABILITY METHODOLOGY}
\label{ch:methods}

\note{Defines every measurement instrument once: activation capture,
linear probing with control baselines, causal ablation, activation
patching, CKA (including the chosen definition of ``matched depth''
across heterogeneous models), and training-dynamics tracking.}

\section{Activation Capture}
\section{Linear Probing}
\subsection{Probe Family and Metrics}
\subsection{Control Baselines}
\section{Causal Ablation}
\section{Activation Patching}
\section{Representational Similarity}
\section{Training-Dynamics Tracking}
\section{Compute and Reproducibility}

\chapter{Results I: What Is Encoded}
\label{ch:results-encoding}

\note{Answers RQ1. Layer-by-layer physics encoding profiles for every
architecture and ground-truth quantity, always shown against the control
baselines; compares geometry (distance, direction) with law (force
magnitude, potential).}

\section{Physics Encoding Profiles}
\section{Geometry Versus Law}
\section{Depth of Emergence}
\section{The Hamiltonian Network Versus the Rest}
\section{Unconstrained Baselines}
\section{Summary of RQ1 Findings}

\chapter{RESULTS II: WHAT IS USED}
\label{ch:results-causal}

\note{Answers RQ2. Ablation of probed directions and activation patching
of a charge flip, separating quantities the models merely encode from
quantities they causally use.}

\section{Ablation Results}
\section{Correlate Versus Cause}
\section{Activation Patching: Tracing a Charge Flip}
\section{Encodes Versus Uses}
\section{Summary of RQ2 Findings}

\chapter{RESULTS III: REPRESENTATIONAL GEOMETRY AND DYNAMICS}
\label{ch:results-geometry}

\note{Answers RQ3 and RQ4 with cross-architecture CKA (plus a sensitivity
check on the matched-depth definition), and RQ5 with probe accuracy over
training. RQ5 is exploratory and may be folded into the Discussion.}

\section{Cross-Architecture Similarity}
\section{Does the Implementation of Symmetry Matter?}
\section{Physics-Explicit Versus Symmetry-Constrained}
\section{Sensitivity to the Matched-Depth Definition}
\section{Generalization Dynamics}
\section{Summary of RQ3--RQ5 Findings}

\chapter{DISCUSSION}
\label{ch:discussion}

\note{Revisits the hypothesis against the results, draws implications for
physics ML (representational alignment as a third evaluation axis) and for
interpretability (what analytic ground truth buys), and states
limitations.}

\section{Revisiting the Central Hypothesis}
\section{Implications for Physics-Informed Machine Learning}
\section{Implications for Interpretability}
\section{Learned Law Versus Shortcut, Revisited}
\section{Limitations}

\chapter{Conclusion and Future Work}
\label{ch:conclusion}

\note{Restates the contributions with results in hand and lists future
directions (other force laws, nonlinear probes, symbolic distillation,
larger systems, transfer to molecular models).}

\section{Summary of Contributions}
\section{Future Work}


% ---------------------------- Back matter ----------------------------
\backmatter
\buheading{APPENDIX A: SIMULATOR AND VERIFICATION DETAILS}
\label{app:simulator}

\note{Integrator derivation, softening choice, full verification tables.}
   % appendices BEFORE bibliography
\buheading{APPENDIX B: HYPERPARAMETERS AND TRAINING CURVES}
\label{app:hyperparameters}

\note{Per-model hyperparameter tables and all training curves.}

\buheading{APPENDIX C: PROBE CONTROL BASELINES AND EXTENDED RESULTS}
\label{app:probe-controls}

\note{Control-baseline plots that do not fit in Chapter 6; per-seed variance.}

\buheading{APPENDIX D: ADDITIONAL REPRESENTATIONAL SIMILARITY ANALYSES}
\label{app:cka}

\note{CKA under alternative matched-depth definitions; linear vs.\ RBF kernels.}

\buheading{APPENDIX E: CODE AND DATA AVAILABILITY}
\label{app:code}

\note{Repository, software versions, commands to reproduce every figure.}


% Bibliography. unsrtnat = numbered, order of appearance;
% plainnat = alphabetical by author.
\nocite{*}                                % outline stage: remove once chapters cite
\bibliographystyle{unsrtnat}
\bibliography{references}

% ---------------------------------------------------------------------
%  Curriculum Vitae (REQUIRED of all students)
%  You may place a Vita here instead: change the heading to VITA, write
%  it as a third-person narrative, and use double spacing.
% ---------------------------------------------------------------------
\buheading{CURRICULUM VITAE}

\begin{hangingparas}
\note{Required. Full name, a contact address good for 1--2 years,
education, positions, publications. Three to four pages at most.}
\end{hangingparas}
                     % required

\end{document}
